% ---
% title: A medical test
% tags: bayes, total-probability
% ---

\begin{problem}
	In a country 6\% of the people are infected. A new device tests an
	infected person positive with probability $0.96$, and it tests an
	uninfected person positive with probability $0.05$.
	(a) Find the probability that a randomly chosen person tests positive.
	(b) A person tests positive. What is the probability that this person is
	really infected?
\end{problem}

\begin{solution}
	\idea Split the population into two groups: infected and uninfected. Use
	the law of total probability for (a). Then use Bayes' formula for (b). We
	know the chance of a positive test in each group; Bayes' formula turns this
	around and gives the chance of infection for a person who tests positive.

	Let $I$ be the event that the chosen person is infected, and $T$ the event
	that the person tests positive. The data are
	\[ P(I)=0.06,\qquad P(I^c)=1-0.06=0.94,\qquad P(T\mid I)=0.96,\qquad
		P(T\mid I^c)=0.05. \]
	\step{(a)} Every person is either infected or uninfected, and not both. So
	$I$ and $I^c$ split $\Om$ into two parts, and $T$ splits into the two
	mutually exclusive pieces $TI$ and $TI^c$. Add their probabilities, and
	write each one with the multiplication rule $P(TI)=P(T\mid I)P(I)$:
	\[ P(T)=P(T\mid I)P(I)+P(T\mid I^c)P(I^c)
		=0.96\cdot0.06+0.05\cdot0.94=0.0576+0.047=0.1046. \]
	So about 10.5\% of all people test positive.
	\step{(b)} We need $P(I\mid T)$. By the definition of conditional
	probability, $P(I\mid T)=P(IT)/P(T)$. The numerator is the first term of the
	sum in (a), $P(IT)=P(T\mid I)P(I)=0.0576$. The denominator is the whole
	sum from (a). This is Bayes' formula:
	\[ P(I\mid T)=\frac{P(T\mid I)P(I)}{P(T\mid I)P(I)+P(T\mid I^c)P(I^c)}
		=\frac{0.0576}{0.1046}\approx0.5507. \]
	\begin{answer}
		(a) $P(T)=0.1046$ \quad (b) $P(I\mid T)\approx0.5507$
	\end{answer}
\end{solution}
